\section{A bit-test system with mixing parity queries}\label{sec:bit-tests}

Here clauses have three literal positions, with repetitions permitted.
Repeat literals to pad a nonempty shorter clause; replace an empty clause
by two opposing unit clauses on a fresh variable and then pad them.
A formula with no clauses can be replaced by a single tautological clause.
These operations preserve satisfiability and take linear time.

We first construct the finite tests from which the reduction will use
predicates. A test law consists of a distribution of question lists and,
for each sampled list together with its public sampling data, an
acceptance predicate on the answers. Repeated questions are allowed and
count toward the arity.

\begin{theorem}\label{thm:bit-test-system}
For every fixed $\lambda>0$ there are constants $\eta>0$, $B<\infty$,
and a fixed finite list of bounded-arity test laws with the following
properties. From a Boolean $3$CNF formula one can construct, in
polynomial time, a finite uniform space $\cS$ of bit positions and these
test laws on $\cS$.
\begin{enumerate}[label=\textup{(\roman*)}]
\item Every individual question marginal is at most $B$ times uniform
measure on $\cS$. All samplers are polynomially enumerable finite
weighted lists with rational probabilities; every positive seed
probability is at least inverse-polynomial in the input size.
\item If the formula is satisfiable, there is an assignment
$y:\cS\to\F_2$ accepted by every test in every support.
\item If the formula is unsatisfiable, every assignment fails one of
the test laws with probability greater than $\eta$.
\item There is also an exchangeable law of
$(s_1,s_2,s_3,s_4)\in\cS^4$, with uniform marginals, such that the
conditional expectation operator of any coordinate given any other has
norm at most $\lambda$ on mean-zero $L^2(\cS)$. The assignment in
\textup{(ii)} satisfies
\[
 y_{s_1}+y_{s_2}+y_{s_3}+y_{s_4}=0
 \qquad\text{in }\F_2
\]
throughout the support of this law.
\end{enumerate}
All constants, including the degrees of the polynomial size and time
bounds, depend only on $\lambda$. Two independent copies of the law in
\textup{(iv)}, paired coordinatewise, have the same pair-operator bound
on $\cS^2$ and give the corresponding parity identity for the expressions
$y_i+y_j$.
\end{theorem}

\subsection{The matrix walk and the masked parity law}

Pad the variable set to addresses in $\{0,1\}^d$, with $d\geq1$ and
$d=O(\log(\text{input size}+1))$. Dummy addresses impose no clauses.
Choose the fixed square field $\F=\F_q$ of characteristic two and
the pencil $u_1,\ldots,u_m\in\F^d$ from Lemma~\ref{lem:pencil}.
The field will be sufficiently large and the pencil parameter
$\lambda_0$ sufficiently small. Set
\[
 G_0=\Span_{\F}\{u_j u_j^{\mathsf T}:1\leq j\leq m\},
 \qquad G=\operatorname{Mat}_{b\times b}(G_0).
\]
We regard $G$ as a space of $(bd)\times(bd)$ matrices with $d\times d$
blocks. We may take $b=101$; the allocation of its blocks is described
below. Since the pencil contains the standard basis, $G_0$ contains
every diagonal matrix. Also $\dim_{\F}G\leq b^2m=O(d)$.

Let $\mathsf D$ be the law
\begin{equation}\label{eq:rank-one-direction}
 h=(c\otimes u_j)(e\otimes u_j)^{\mathsf T},
 \qquad c,e\ \text{uniform in }\F^b,\quad j\ \text{uniform in }[m],
\end{equation}
with the three choices independent. This law is invariant under
multiplication by a nonzero scalar, and under simultaneous permutations
of block rows and block columns. It may have an atom at zero.
Write
$\mathsf T f(A)=\E_{h\sim\mathsf D}f(A+h)$, and give $G$ uniform measure
$\mu$.

The internal-edge estimate below is the usual centered-indicator
spectral calculation; see Alon and Chung \cite[Lemma~2.3]{AlonChung1988}.
We prove the eigenvalue bound for this particular matrix walk directly.

\begin{lemma}\label{lem:matrix-walk}
On mean-zero $L^2(G)$, the operator $\mathsf T$ has norm at most
$\lambda_1=\lambda_0+q^{-1}$. In particular, for every nonempty
$D\subseteq G$,
\begin{equation}\label{eq:internal-edge-bound}
 \frac{\ip{\ind_D}{\mathsf T\ind_D}}{\mu(D)}
 \leq \mu(D)+\lambda_1.
\end{equation}
\end{lemma}

\begin{proof}
Fix a nonzero $\F_2$-linear functional $\ell:\F\to\F_2$.
Every additive character of $G$ is
$A\mapsto(-1)^{\ell(L(A))}$ for an $\F$-linear functional $L:G\to\F$.
Indeed the pairing $(a,x)\mapsto\ell(ax)$ is nondegenerate: for $a\ne0$,
choose $x$ with $\ell(ax)=1$. Choosing a basis of $G$ and counting then
gives all its additive characters.

For such an $L$, evaluation on Equation~\ref{eq:rank-one-direction}
gives $L(h)=c^{\mathsf T}M_j e$, where every entry of $M_j$ is a
homogeneous quadratic form in $u_j$. If $L\ne0$, some entry is nonzero
on the list: the directions span $G$, as is seen by taking $c,e$ to be
coordinate vectors. That entry vanishes on at most $\lambda_0 m$
indices. For $M_j\ne0$, averaging the character first over $c$ gives
zero unless $M_j e=0$, so its full average is
$q^{-\rank M_j}\leq q^{-1}$. The remaining indices contribute at most
one. Orthogonality of additive characters diagonalizes $\mathsf T$ and
proves the norm bound. Applying it to
$\ind_D-\mu(D)$ proves Equation~\ref{eq:internal-edge-bound}.
\end{proof}

Our field oracle will be a function $P:G\to\F$. Its binary encoding has
positions
\[
 \cS=G\times\F^*,\qquad
 y_{(A,\sigma)}=\ell(\sigma P(A)).
\]
The vector of all $q-1$ masks determines $P(A)$ uniquely by the same
nondegenerate pairing. Any field query below is implemented by querying
all these masks and rejecting unless they are the encoding of a field
element.

Choose a uniformly ordered tuple of four distinct times
$t_1,t_2,t_3,t_4\in\F$, an independent uniform $A\in G$, an independent
$h\sim\mathsf D$, and an independent uniform $\sigma\in\F^*$. Put
\begin{equation}\label{eq:four-masked-positions}
 s_i=(A+t_i h,\sigma\alpha_i),\qquad
 \alpha_i=\left(\prod_{j\ne i}(t_i-t_j)\right)^{-1}.
\end{equation}
This law is exchangeable and each marginal is uniform. If the restriction
of $P$ to every sampled line is a polynomial of degree at most two,
Lagrange interpolation gives
$\sum_i\alpha_iP(A+t_i h)=0$; applying $\ell(\sigma\,\cdot)$ gives the
required binary parity identity.

\begin{lemma}\label{lem:masked-pair-mixing}
Every pair-operator of Equation~\ref{eq:four-masked-positions} has
mean-zero norm at most
\[
 \max\{\lambda_1,\,2/(q-3)\}.
\]
\end{lemma}

\begin{proof}
By exchangeability it suffices to consider the first two times,
denoted $a,b$, and write the other two as $c,d$. For every fixed
$(a,b,c,d,h)$ and every $(x,m)\in G\times\F^*$,
\[
 \Prob[(A+ah,\sigma\alpha_a)=(x,m)\mid a,b,c,d,h]
 =\frac1{|G|(q-1)}.
\]
Conditioning on the complete first bit position therefore preserves
the joint time and direction law. The second position is
\[
 \bigl(x+(b-a)h,\ m\alpha_b/\alpha_a\bigr).
\]
For every time tuple, $(b-a)h$ has law $\mathsf D$. Consequently its
law is independent of the mask multiplier, and the pair operator is
the tensor product of $\mathsf T$ and the averaged mask permutation.

For fixed distinct $a,b,c$, let $k=(b-c)/(a-c)$. Direct cancellation,
using characteristic two, gives
\[
 \frac{\alpha_a}{\alpha_b}
 =k\frac{b-d}{a-d}.
\]
The fractional-linear map on the right sends the excluded projective
arguments $a,b,c,\infty$ to $\infty,0,k^2,k$, respectively. It is
one-to-one on the projective line. Thus, for allowed $d$, this ratio is
uniform on $\F^*\setminus\{k,k^2\}$. Here $k\notin\{0,1\}$, so the two
omitted values are distinct. The inverse ratio likewise omits exactly
two nonzero values.

On mean-zero functions of the uniform mask, the sum of all $q-1$
multiplicative permutations is zero. The average with two permutations
removed therefore has norm at most $2/(q-3)$, since every permutation is
an $L^2$ isometry. Averaging over $a,b,c$ preserves this bound. The
orthogonal subspaces of functions spatially centered at each mask and
of mask-only mean-zero functions are invariant under the tensor
operator. The former has norm at most $\lambda_1$, and the latter has
the bound just proved. These subspaces exhaust the orthogonal
complement of the constants.
\end{proof}

Choose the field and pencil so that
\begin{equation}\label{eq:fixed-field-smallness}
 \lambda_1<10^{-6},\qquad
 \max\{\lambda_1,2/(q-3)\}\leq\lambda,\qquad
 \frac6{q-1}\leq\frac1{16}.
\end{equation}
These are fixed choices. On two independent copies of the bit space,
the pair operator is the tensor square. Decomposing each factor into
constants and mean-zero functions shows that its nonconstant norm is
at most the same bound. This proves the final assertion of
Theorem~\ref{thm:bit-test-system} once the promised line-polynomial
assignment is constructed.

\subsection{Boolean tables and polynomial slack}

Evaluation codes formed from spans of matrix minors were studied by
Beelen, Ghorpade, and H\o holdt as affine Grassmann codes
\cite{BeelenGhorpadeHoholdt2010}. Here we use principal minors to extend
Boolean tables, and products of two minors to retain quadratic restrictions
on the chosen rank-one lines.

Let $T$ be the principal matrix on the first three blocks, and put
$n=3d$. For $I\subseteq[n]$, define
$p_I(T)=\det T_{I,I}$, including $p_\varnothing=1$.
On a Boolean diagonal $T=\diag(z)$ these are the monomials
$\prod_{i\in I}z_i$. Every table on $\{0,1\}^n$ has a unique expansion
in these monomials: at the indicator of a set $J$, its value is
$\sum_{I\subseteq J}a_I$, which determines the coefficients successively
by increasing $|J|$. We extend a Boolean table to matrices by replacing
its monomials by the corresponding principal minors.

For each $a=(a_1,a_2,a_3)\in\{0,1\}^3$, let $E_a(T)$ be the public
extension of the following table on triples of addresses. It equals
one if the input contains a clause on that ordered triple whose
falsifying answer pattern is $a$, and zero otherwise. Multiple clauses
and repeated variables cause no change to this definition.

There is one assignment field $X$, required to depend only on the first
diagonal block. For $i=1,2,3$, let $\pi_i$ simultaneously swap block
one with block $i$ in rows and columns, with $\pi_1$ the identity, and
write $X_i(A)=X(\pi_i A)$. For each $a$ define
$M_{a,i}=X_i$ when $a_i=1$ and $M_{a,i}=1-X_i$ when $a_i=0$.
Introduce gate fields $Y_a,Z_a$. On Boolean diagonal $T$ we want
\begin{equation}\label{eq:clause-identities}
 X^2-X=0,\qquad
 Y_a-E_aM_{a,1}=0,\qquad
 Z_a-M_{a,2}M_{a,3}=0,\qquad Y_aZ_a=0.
\end{equation}
There are $25$ identities. In a satisfying assignment, extend the
assignment table to $X$ using the first block. Extend the Boolean
tables $E_aM_{a,1}$ and $M_{a,2}M_{a,3}$ to $Y_a,Z_a$ using $T$.
Every left side of Equation~\ref{eq:clause-identities} is then a linear
combination of products $p_Ip_J$, and vanishes on all Boolean diagonals.
For the last identity, simultaneous nonzero values would exhibit the
forbidden pattern of a clause.

We need to enforce these identities at uniform matrix arguments.
The following explicit reduction supplies their slack.

\begin{lemma}\label{lem:minor-product-slack}
A linear combination of products $p_I(T)p_J(T)$ that vanishes on all
Boolean diagonal matrices is a sum of terms of the form
\begin{equation}\label{eq:slack-coefficient-patterns}
 cMN\,T_{ij},\qquad
 cMN\,T_{ij}(T_{ii}-1),\qquad
 cMN\,(T_{ii}^2-T_{ii}),
\end{equation}
where $i\ne j$ in the first two forms, $c\in\F$, and $M,N$ are
arbitrary square minors of $T$, possibly empty.
\end{lemma}

\begin{proof}
For $i\in K$, write the row-$i$ expansion
\[
 p_K=T_{ii}p_{K\setminus\{i\}}+
       \sum_{j\ne i}T_{ij}C^K_{ij},
\]
where $C^K_{ij}$ is the signed cofactor when $j\in K$, and zero
otherwise. Every nonzero cofactor is a scalar sign times an arbitrary
minor. We retain minus signs in the identities below although the
field has characteristic two.

If $i\in J\setminus I$, expansion of $p_J$ and $p_{I\cup\{i\}}$
along row $i$ cancels the diagonal terms and gives
\begin{align}
 p_Ip_J-p_{I\cup\{i\}}p_{J\setminus\{i\}}
 =\sum_{j\ne i}T_{ij}
 \bigl(p_I C^J_{ij}
       -C^{I\cup\{i\}}_{ij}p_{J\setminus\{i\}}\bigr).
 \label{eq:minor-move}
\end{align}
If $i\in I\cap J$, first expand $p_J-p_{J\setminus\{i\}}$ and then
expand the $p_I$ multiplying $T_{ii}-1$. This gives
\begin{align}
 p_Ip_J-p_Ip_{J\setminus\{i\}}
 ={}&(T_{ii}^2-T_{ii})
       p_{I\setminus\{i\}}p_{J\setminus\{i\}}\notag\\
 &+\sum_{j\ne i}T_{ij}p_I C^J_{ij}
 +\sum_{j\ne i}T_{ij}(T_{ii}-1)
       C^I_{ij}p_{J\setminus\{i\}}.
 \label{eq:minor-delete}
\end{align}
Each residual term has a form in
Equation~\ref{eq:slack-coefficient-patterns}.

Starting from $(I,J)$, choose an element of the second set and apply
Equation~\ref{eq:minor-move} or Equation~\ref{eq:minor-delete}.
The new second set loses that element, and the union of the sets is
unchanged. After exactly $|J|$ steps, the remaining product is
$p_{I\cup J}$. Telescoping therefore reduces every original product to
this canonical principal minor plus permitted residuals.

Applying the reduction to the given linear combination leaves
$\sum_U b_U p_U$ plus permitted residuals. All residuals vanish on
Boolean diagonals. At the indicator diagonal of a set $S$ we obtain
$\sum_{U\subseteq S}b_U=0$. Induction on $|S|$, beginning with the
empty set, gives $b_S=0$ for every $S$. Thus all canonical terms
cancel.
\end{proof}

Reserve two disjoint tag banks, each consisting of three new blocks.
Let $b_r(i)$ be the copy of address index $i$ in bank $r\in\{1,2\}$.
For an arbitrary minor $M(T)=\det T_{R,C}$, define
\[
 \mathsf E^{(r)}_{ij}M(A)
 =\det A_{(R,b_r(i)),(C,b_r(j))},
\]
where the new row and column are appended to their respective lists.
For the empty minor this is just the appended entry. If $V$ is
supported in the principal bank $r$, the extended submatrix contains
exactly one entry of that bank, its appended corner. Its corner
cofactor is $M$. Consequently, with
$\Delta_VL(A)=L(A+V)-L(A)$,
\begin{equation}\label{eq:tag-extraction}
 \Delta_V(\mathsf E^{(r)}_{ij}M)
 =V_{b_r(i),b_r(j)}M.
\end{equation}

For each of the $25$ identity differences, group the residual terms
from Lemma~\ref{lem:minor-product-slack} by the three forms in
Equation~\ref{eq:slack-coefficient-patterns}. Construct three fields
$L_1,L_2,L_3$ as follows, summing with the same coefficients as the
residuals:
\[
 \begin{array}{c|c}
 \text{residual term}&\text{term placed in the field}\\ \hline
 cT_{ij}MN&c(\mathsf E^{(1)}_{ij}M)N\quad\text{in }L_1\\
 cT_{ij}(T_{ii}-1)MN&
 c(\mathsf E^{(1)}_{ij}M)(\mathsf E^{(2)}_{ii}N)
       \quad\text{in }L_2\\
 c(T_{ii}^2-T_{ii})MN&
 c(\mathsf E^{(1)}_{ii}M)N\quad\text{in }L_3.
 \end{array}
\]
Let $v(T)$ copy the off-diagonal part of $T$ into bank one,
$w(T)$ copy $\diag(T)-I$ into bank two, and $u(T)$ copy
$\diag(T)^2-\diag(T)$ into bank one. Here $\diag(T)$ means the
diagonal matrix retaining the diagonal entries of $T$. These shifts
belong to $G$: copying a full block preserves membership in $G_0$, and
removing or replacing diagonal entries is allowed because every
diagonal matrix belongs to $G_0$.
All these shifts leave $T$ unchanged. Equation~\ref{eq:tag-extraction}
therefore proves the identity
\begin{equation}\label{eq:identity-with-slack}
 F(A)=\Delta_{v(T)}L_1(A)
       +\Delta_{v(T)}\Delta_{w(T)}L_2(A)
       +\Delta_{u(T)}L_3(A)
\end{equation}
for the honest difference $F$. In addition, $L_1,L_3$ are affine
under all translations in bank one, and $L_2$ is separately affine
under translations in each bank. Indeed its first extended factor
depends affinely on bank one and not on bank two, and its second
factor has the reverse properties. Cross entries between a tag bank
and the address indices do not change under principal-bank
translations.

\subsection{One oracle, a fixed list of tests, and completeness}

The one assignment field, sixteen gate fields, and seventy-five slack
fields are $92$ fields altogether. Give each its own selector index,
using one index in each of $92$ fresh blocks. Together with the three
address and six tag blocks, this uses $101$ blocks.
Every honest field constructed above is a linear combination of
products of two arbitrary minors using address and tag indices only.
For a field $F$ with selector index $s_F$, replace the first minor in
each of its terms $cMN$ by the minor with appended row and column
$s_F$. Sum all such packed terms over the $92$ fields to obtain $P$.
If $v_F$ is the unit diagonal matrix at $s_F$, the corner-cofactor
calculation gives
\begin{equation}\label{eq:selector-isolation}
 \Delta_{v_F}P=F.
\end{equation}
Terms belonging to another field omit $s_F$ completely, so they make
no contribution to this difference. The shift $v_F$ belongs to $G$.

Every minor of $A+th$, with $h$ as in
Equation~\ref{eq:rank-one-direction}, is affine in $t$. To see this,
expand the determinant by multilinearity in columns. Terms using two
perturbed columns vanish because those columns are proportional.
Thus each packed product has line degree at most two, as does $P$.

For an arbitrary oracle $P$, use Equation~\ref{eq:selector-isolation}
as the \emph{definition} of all proof fields. The ordinary laws are:
\begin{enumerate}
\item The line law samples uniform $A$ and $h\sim\mathsf D$, queries
$P(A+th)$ for all $t\in\F$, and checks degree at most two.
\item Each of the $25$ laws checks
Equation~\ref{eq:identity-with-slack} at uniform $A$, with its own
three designated slack fields and the corresponding difference from
Equation~\ref{eq:clause-identities}.
\item The assignment invariance law samples a uniform translation
$v\in G$ whose first diagonal block is zero, and checks
$X(A+v)=X(A)$ at uniform $A$.
\item For each slack field and each bank in which it is required to
be affine, let $B_r$ be the subspace of $G$ supported in that
principal bank. Test
\begin{equation}\label{eq:bank-tests}
 \Delta_w\Delta_vL(A)=0
 \quad(v,w\text{ independent uniform in }B_r),
 \qquad
 \Delta_{av}L(A)=a\Delta_vL(A)
 \quad(v\text{ uniform in }B_r)
\end{equation}
at uniform $A$, with a separate law for each fixed scalar $a\in\F$.
\end{enumerate}
Every displayed field value is evaluated through its selector
difference and all masks. Queries are made even when a coefficient
vanishes, keeping each law's list length fixed.

All master-oracle query points have uniform marginals on $G$.
Translations and block swaps are bijections. The maps
$A\mapsto A+v(T(A))$, and the analogous tag shears, are bijections
because they leave their source $T$ fixed; subtraction of the same
shift is their inverse. Their compositions with selectors have the
same property. Bank-test parameters are independent of the uniform
base. A fixed mask at a uniform base is at most $q-1$ times uniform
on $\cS$, so we may use $B=q-1$.

The honest construction above satisfies every law identically and
has line degree at most two. It therefore proves perfect completeness,
including the auxiliary parity identity. The witness-dependent fields
are used only to prove completeness; all public test laws are specified
without knowing a satisfying assignment.

\subsection{Deterministic correction of almost quadratic lines}

The remaining issue is constant soundness. It is not enough for the
identities to hold at most matrices, since the Boolean diagonals may
be sparse. We first obtain exact line identities by correcting a small
set. This is a local-to-global correction problem in the self-testing
framework of Blum, Luby, and Rubinfeld \cite{BlumLubyRubinfeld1993}.
The modal interpolation and transverse-grid consistency argument in the
point-filling step has precedents in polynomial self-testing
\cite{GemmellEtAl1991,RubinfeldSudan1996}. The next lemma proves the
required correction statement for our restricted distribution of
directions, with one common high-mass set of good directions.

\begin{lemma}\label{lem:quadratic-correction}
Put $\gamma=10^{-3}$. Suppose the line law of a function $P:G\to\F$
fails with probability at most $\delta^2$, where $\delta>0$ satisfies
\begin{align}
 2\delta+\lambda_1&<\gamma/2,\notag\\
 1-4\gamma-4\delta-2/(q-1)&>\delta+\lambda_1,\label{eq:correction-smallness}\\
 27\gamma+5\delta&<1.\notag
\end{align}
There is a scalar-invariant set $H$ of directions of
$\mathsf D$-mass at least $1-\delta$, and a function $P'$ differing
from $P$ on at most $2\delta$ of $G$, such that
$t\mapsto P'(A+th)$ has degree at most two for every $A$ and $h\in H$.
\end{lemma}

\begin{proof}
Let $H$ consist of directions whose fraction of failing bases is at
most $\delta$. Mark initially those points whose fraction of failing
directions is greater than $\delta$, obtaining $D_0$ with
$\mu(D_0)\leq\delta$. Markov's inequality also gives
$\mathsf D(H)\geq1-\delta$. Both notions of line failure are unchanged
by multiplying a direction by a nonzero scalar.

Enlarge $D_0$ by adding an outside point $x$ whenever
$\Prob_h[x+h\in D]\geq\gamma$. If $D_0$ has $k\geq1$ points, the
first set of size $2k$, if reached, has internal directed edge mass
at least $k\gamma/|G|$: every added source point contributed that
amount at its addition and those edges remain internal. Dividing by
its density gives at least $\gamma/2$, contrary to
Equation~\ref{eq:internal-edge-bound} and
$2\delta+\lambda_1<\gamma/2$. If $k=0$, no point can be added.
Thus the final $D$ has density at most $2\delta$, and every point
outside it has entry probability less than $\gamma$ into $D$.
Call these outside points clean.

Fix $h\in H$ and four distinct times $t_i$. Let $J$ be the set of
bases whose four positions are clean but for which
$\sum_i\alpha_iP(A+t_i h)\ne0$, using the interpolation coefficients
of Equation~\ref{eq:four-masked-positions}. Then $\mu(J)\leq\delta$,
because every such base has a failing $h$-line. For $A\in J$, put
$z_i=A+t_i h$. Sample an independent $k\sim\mathsf D$ and a uniform
$a\in\F^*$. Since $z_i\notin D_0$, all four functions
$s\mapsto P(z_i+sk)$ are quadratic except on an event of probability
at most $4\delta$. On that event's complement their relation
discrepancy is a degree-two polynomial in $s$, nonzero at zero, and
so is zero at at most two choices of $a\in\F^*$.
Separately, each $z_i+ak$ is missing with probability at most
$\gamma$, since $ak$ has exactly law $\mathsf D$. We do not condition
these cleanliness bounds on line goodness. A union bound gives
\[
 \Prob_{k,a}[A+ak\in J]\geq
 1-4\delta-4\gamma-\frac2{q-1}.
\]
The step $ak$ has law $\mathsf D$, including at zero, so a nonempty
$J$ would contradict Equation~\ref{eq:internal-edge-bound} and
Equation~\ref{eq:correction-smallness}. All four-position relations
whose positions are clean are therefore exact.

We now fill the missing points one at a time. Maintain that every
assigned point has entry probability at most $\gamma$ into the
current missing set and that every $H$-direction four-position
relation on assigned points is exact. If the missing set is
nonempty, Equation~\ref{eq:internal-edge-bound} supplies a point $x$
inside it whose entry probability is at most $\gamma$. Removing $x$
can only decrease all entry probabilities.

Fix three distinct nonzero times $a_1,a_2,a_3$ and interpolation
coefficients $\ell_1,\ell_2,\ell_3$ for evaluation at zero. A direction
$k$ predicts $\sum_i\ell_iP'(x+a_i k)$ if its three positions are
already assigned, and gives an arbitrary fixed default otherwise.
For independent $h,k$, require the six axis points
$x+a_i h,x+a_j k$ and the nine grid points
$x+a_i h+a_j k$ to be assigned. The six axis failures cost at most
$6\gamma$. After an assigned first-axis point is exposed, the
independent second direction reaches the missing set with probability
at most $\gamma$; the nine additional failures cost at most
$9\gamma$. No independence among these fifteen events is asserted.
Requiring $h,k\in H$ costs at most $2\delta$.
On the resulting event, interpolation of the grid first by rows and
then by columns shows that the two predictions agree, because all
the relations used involve already assigned points.

Let $\pi_x$ be the full prediction distribution, defaults included.
Its collision probability is at least $1-15\gamma-2\delta$ and is
at most its largest atom. Choose one modal value $m_x$ once and
assign it to $x$. A random prediction equals $m_x$ and uses an
already assigned triple except with probability at most
$18\gamma+2\delta$. At any old assigned center, the analogous
prediction reproduces its value with an old assigned triple except
with probability at most $3\gamma+\delta$.

Consider any fixed relation newly enabled by assigning $x$. A zero
direction gives an automatic constant-line identity. Otherwise its
four positions are distinct and only one is $x$. For a fresh
transverse direction $k$, simultaneous correct interpolation from
old assigned triples at its four centers fails with probability at
most
\[
 (18\gamma+2\delta)+3(3\gamma+\delta)=27\gamma+5\delta<1.
\]
There is therefore such a $k$. Each of the three translated copies
of the relation lies entirely in the old assigned set and is zero.
Interpolating those three identities back recovers the old values
and the already selected $m_x$, so the original relation is zero.
Repeated grid positions cause no problem: they denote the same
assigned value. A grid position equal to the still-missing center
was already excluded by the old-assignment event.

The modal value is fixed before the relation is considered. Different
transverse directions may establish different deterministic
relations for that same value; no union bound over relations or
peeling steps is used. This proves the induction. At its completion
all four-position relations hold in every direction in $H$.
Interpolation at three distinct parameters and then at each fourth
parameter proves the asserted line degree. Only the original enlarged
set, of density at most $2\delta$, was changed.
\end{proof}

\subsection{Support gaps and exact bank affinity}

\begin{lemma}\label{lem:line-support-gap}
Suppose a scalar-invariant set of directions has
$\mathsf D$-mass at least $1-\rho$, and $F:G\to\F$ has degree at most
$D<q-1$ on every line in those directions. Then either $F=0$
identically or
\[
 \mu(\supp F)\geq 1-\rho-\frac{D}{q-1}-\lambda_1.
\]
\end{lemma}

\begin{proof}
From a point $x$ of the support, sample $h\sim\mathsf D$ and an
independent uniform $a\in\F^*$. On a good direction, the nonzero
line polynomial has at most $D$ roots among those scalars. Thus
$x+ah$ is in the support with probability at least
$1-\rho-D/(q-1)$. The law of $ah$ is the original direction law.
Average over support points and apply
Equation~\ref{eq:internal-edge-bound}.
\end{proof}

Suppose now that every ordinary binary test law fails with probability
at most $\eta$. Decode the full mask vector to $P(A)$ wherever it is
valid, and choose $P(A)$ arbitrarily otherwise. Acceptance of a binary
test certifies both valid encodings at all its queried points and its
field relation. Thus every decoded field law also fails with
probability at most $\eta$.

Let $Q$ be a fixed upper bound on the number of master-field queries
in a law, counting repetitions; for example $Q=q+32$ suffices for
the displayed list. Put $\delta=\sqrt\eta$ and apply
Lemma~\ref{lem:quadratic-correction}. Replacing $P$ by $P'$ increases
any field-law failure probability by at most $2Q\delta$, because
each master query is uniform. Write
\begin{equation}\label{eq:corrected-test-error}
 \tau=\eta+2Q\delta.
\end{equation}
Intersect the pullbacks of $H$ under the three block permutations
$\pi_1,\pi_2,\pi_3$. Their common scalar-invariant direction set has
missing mass at most $K_{\rm sw}\delta$, with $K_{\rm sw}=3$.
Every fixed translate of every corrected proof field, and of every
$X_i$, has line degree at most two on this common set.
Lemma~\ref{lem:line-support-gap}, with the conservative degree bound
six, gives the support lower bound
\begin{equation}\label{eq:kappa-support}
 \kappa=1-K_{\rm sw}\delta-\frac6{q-1}-\lambda_1
\end{equation}
for any nonzero discrepancy to which that degree bound applies.

First consider a fixed parameter in the assignment invariance test.
Its discrepancy has degree at most two. If it is not identically
zero, it fails on at least a $\kappa$ fraction of bases. Hence at
least a $1-\tau/\kappa$ fraction of the tested translations preserve
$X$ identically. Those translations form an additive subgroup.
A subgroup of a finite additive group with density greater than
one half is the whole group, since otherwise one of its disjoint
cosets has the same size. Thus the assignment invariance is exact
provided $\tau/\kappa<1/2$.

For a bank $B_r$ and slack field $L$, fix $v$ and let
\[
 W_v=\{w\in B_r:
       \Delta_vL(A+w)=\Delta_vL(A)\text{ for every }A\}.
\]
This is the translation stabilizer of the function $\Delta_vL$,
hence a subgroup. Each fixed $(v,w)$ discrepancy has line degree
at most two. The fraction of pairs for which it is not identically
zero is at most $\tau/\kappa$. If $\tau/\kappa<1/16$, Markov's
inequality leaves more than three quarters of $v$ for which more
than three quarters of $w$ are valid. For these $v$, $W_v=B_r$.
The set
\[
 V=\{v\in B_r:\Delta_vL\text{ is invariant under all }B_r
                    \text{ translations}\}
\]
is also a subgroup. Indeed, for $u,v\in V$, the difference cocycle
gives
\[
 \Delta_{u+v}L(A)
 =\Delta_uL(A+v)+\Delta_vL(A)
 =\Delta_uL(A)+\Delta_vL(A),
\]
and $\Delta_{-v}L(A)=-\Delta_vL(A-v)=-\Delta_vL(A)$.
Thus $V=B_r$, and the map $v\mapsto\Delta_vL(A)$ is additive for
every $A$.

For each fixed scalar $a$, the map
$v\mapsto(\Delta_{av}L(A)-a\Delta_vL(A))_{A\in G}$ is now additive.
Its kernel has density at least $1-\tau/\kappa>1/2$ by that scalar's
test and the support gap, and so equals $B_r$. This proves exact
$\F$-affinity in each required bank. There are only $q$ scalar laws,
with $q$ fixed. The argument has exactified whole translation laws
before introducing a basis of any bank.

We may now expand the shifts in
Equation~\ref{eq:identity-with-slack}. For bases $(e_i)$ and $(f_j)$
of the two banks, exact separate affinity gives
\begin{equation}\label{eq:exact-bilinear-shift}
 \Delta_{v(T)}\Delta_{w(T)}L(A)
 =\sum_{i,j}v_i(T)w_j(T)\Delta_{e_i}\Delta_{f_j}L(A).
\end{equation}
Each coefficient function on the right is a fixed combination of
translates of a selector difference of $P'$, so has line degree at
most two. The coordinates of $v(T)$ and $w(T)$ have degree at most
one on an $A$-line. The corresponding coordinates of $u(T)$ have
degree at most two. The one-bank expansions have the same form
with one fixed difference. Thus every slack term has line degree
at most four. The main field products have degree at most four,
and every public principal minor has degree at most one.
The number of coefficients in Equation~\ref{eq:exact-bilinear-shift}
does not affect the degree of their sum; every coefficient identity
is already exact.

Every full identity discrepancy therefore has degree at most four
on the common good directions and fails on at most a $\tau$ fraction
of bases. If $\tau<\kappa$, the support gap forces it to vanish
identically. On Boolean diagonal $T$, the shifts $v(T)$ and $u(T)$
are zero. Both single differences and the mixed difference with
$v(T)=0$ therefore vanish, regardless of $w(T)$, and
Equation~\ref{eq:clause-identities} holds exactly.
Every triple of Boolean addresses is available because $G_0$
contains all diagonal matrices. Exact invariance makes $X$ a
well-defined function of its first block alone. The Booleanity
identity assigns a bit to each Boolean address, and the gate
identities exclude every falsifying clause pattern: if all three
assigned bits matched a forbidden pattern, then $E_a=M_{a,1}
=M_{a,2}=M_{a,3}=1$, forcing $Y_a=Z_a=1$ and contradicting
$Y_aZ_a=0$. This is a satisfying assignment to the input formula.

All required smallness conditions follow from a single fixed choice.
With Equation~\ref{eq:fixed-field-smallness} in force, put
\[
 d_0=\min\left\{\frac{\gamma}{16},
          \frac1{16K_{\rm sw}},\frac1{128(1+2Q)}\right\},
 \qquad 0<\eta\leq d_0^2.
\]
Then Equation~\ref{eq:correction-smallness} holds,
$\kappa\geq1-1/16-1/16-\lambda_1>1/2$, and
$\tau\leq(1+2Q)\sqrt\eta\leq1/128$.
In particular $\tau/\kappa<1/16$ and $\tau<\kappa$.
No bound contains $d$ or $|G|$. An unsatisfiable formula must
therefore have a law failing with probability greater than this
$\eta$.

Finally, the space $G$ has size $q^{O(d)}=2^{O(d)}$, and all
test samplers enumerate fixed products of this space, fixed finite
fields, and the $O(d)$-length pencil. Public table interpolation
uses at most $2^{3d}$ addresses and elementary linear algebra.
Every arithmetic operation has polynomial bit complexity; seed
probabilities are rational products of the reciprocals of these
polynomial-size spaces. Query arity, number of laws, and all mask
counts are constants. The pencil construction itself takes
$2^{O(d)}$ time by Lemma~\ref{lem:pencil}.
These observations prove the size and sampler assertions and
complete the proof of Theorem~\ref{thm:bit-test-system}.
